\section{Introduction}

\subsection{State of the art}
    Chemical propulsion offers the significant advantage of generating high thrust levels, which are essential for rocket launch and the insertion of spacecraft into orbit. However, once in orbit, it is legitimate to question whether these propulsion methods remain the most efficient option.

The energy delivered by chemical combustion is intrinsically linked to the mass of propellant available on board. This dependence introduces a major limitation: increasing the available energy requires additional propellant mass, yet accelerating this extra mass demands even more energy, leading to a compounding cycle.

As a result, this strong mass dependence fundamentally constrains the performance of chemical propulsion systems, particularly their specific impulse, which is typically limited to an upper bound on the order of 500 seconds (\cite{mtp_book} p.652).

We will see that this limit can be overcome thanks to electric propulsion.

\subsection{Scope and Objectives}
The purpose of this report is to study and compare five existing electric propulsion systems currently used in industry. These technologies are based on three fundamental physical principles: electrostatic, electrothermal, and electromagnetic propulsion (\cite{mtp_book} p.651).

For each technology, this report presents the underlying physical principles, its practical implementation, and key performance parameters such as thrust ($F$), specific impulse ($I_{sp}$), exhaust velocity ($u_e$), and efficiency ($\eta$). Finally, a comparative analysis is conducted to highlight the advantages, limitations, and typical applications of each technology.